\chapter{Toxoplasmosis}
\index{Toxoplasmosis}
\medskip
\noindent\textit{Educational reference; always seek professional advice for individual care.}

\section*{Definition and Overview}
Toxoplasmosis is an infection caused by the parasite \textit{Toxoplasma gondii}. Many people have no symptoms or only a mild flu-like illness. The infection can be serious in pregnancy, in a fetus or newborn, in people with a weakened immune system, and in the eye. After infection, the parasite can remain dormant in body tissues and reactivate when immunity is severely reduced.

\section*{How It Spreads}
People can acquire toxoplasmosis by eating undercooked meat containing tissue cysts, swallowing microscopic material from contaminated soil or cat faeces, eating unwashed produce, or drinking contaminated water. Infection can pass from a pregnant person to the fetus when infection occurs during or shortly before pregnancy. Rare routes include organ transplantation and blood transfusion. Cats are part of the parasite's life cycle, but touching or living with a cat does not by itself mean infection.

\section*{Clinical Features}
\begin{itemize}
    \item No symptoms, or fever, tiredness, headache, muscle aches, and swollen lymph nodes
    \item Blurred vision, eye pain, floaters, or light sensitivity from ocular disease
    \item Severe headache, confusion, weakness, seizures, or reduced consciousness in advanced immunosuppression
    \item In congenital infection: eye disease, enlarged liver or spleen, jaundice, seizures, or developmental problems; some effects appear later
\end{itemize}

\section*{When to Seek Medical Care}
Seek medical advice for persistent swollen glands or flu-like symptoms if pregnant or immunocompromised, and for any new visual symptom. \textbf{Contact local emergency services} for seizure, confusion, severe headache with weakness, reduced consciousness, or sudden major visual loss.

\section*{Diagnosis}
Diagnosis depends on the clinical situation. Blood tests may look for antibodies, but results can require repeat testing or specialist interpretation. Eye disease is diagnosed by an eye examination. Brain imaging and other tests may be needed in people with severe immune suppression or neurological symptoms. In pregnancy, specialist assessment and fetal imaging may be needed; a positive screening result does not by itself prove fetal infection.

\section*{Management}
Healthy people with mild illness often need rest, fluids, and observation rather than antiparasitic treatment. Persistent, ocular, congenital, or severe disease may require combinations such as pyrimethamine, sulfadiazine, and folinic acid under specialist supervision. Treatment in pregnancy is individualised according to timing, testing, and fetal findings. People with HIV or another cause of immunosuppression also need management of the underlying condition. Do not take antiparasitic medicines without medical guidance.

\section*{Prevention}
\begin{itemize}
    \item Cook meat thoroughly and avoid tasting meat before it is fully cooked
    \item Wash hands, utensils, and surfaces after handling raw meat; wash fruit and vegetables
    \item Wear gloves while gardening and wash hands afterwards
    \item If pregnant or immunocompromised, avoid cleaning a cat-litter tray if possible; otherwise use gloves, wash hands, and change litter daily
    \item Keep cats indoors where possible and do not feed them raw meat
    \item Follow local advice about safe drinking water
\end{itemize}

\section*{Outlook}
Most immunocompetent people recover fully. Congenital, ocular, and central-nervous-system disease can cause lasting complications, but early diagnosis and specialist treatment improve outcomes. Recurrence is possible when immunity becomes severely weakened.

\section*{Sources and Further Reading}
\begin{itemize}
    \item Centers for Disease Control and Prevention. \textit{Toxoplasmosis.} https://www.cdc.gov/toxoplasmosis/
    \item World Health Organization. \textit{Toxoplasmosis information.} https://www.who.int/
    \item NHS. \textit{Toxoplasmosis.} https://www.nhs.uk/conditions/toxoplasmosis/
    \item Merck Manual. \textit{Toxoplasmosis.} https://www.merckmanuals.com/
\end{itemize}
